\subsection{Implementácia mapového modulu a geolokalizácie}

Ako základ pre mapu je použitá knižnica \textit{Google Maps SDK}, ktorá poskytuje všetky potrebné funkcionality pre prácu s geografickými dátami. Používateľ na hlavnej obrazovke vidí interaktívnu mapu, na ktorej sú jednotlivé udalosti zobrazené vo forme markerov.

\subsubsection{Zobrazenie mapy a integrácia s Compose}
Keďže \texttt{MapView} nie je natívny komponent prostredia \textit{Jetpack Compose}, v aplikácii je využitý \texttt{AndroidView}, ktorý slúži ako most (\textit{bridge}) medzi klasickým \textit{Android View} systémom a \textit{Compose UI}. Mapa sa pri prvotnom načítaní predvolene centruje na súradnice Bratislavy.

\begin{lstlisting}[language=Kotlin, caption=Inicializácia MapView v prostredí Compose]
val mapView = remember { MapView(context).also { it.onCreate(null) } }

AndroidView(factory = { mapView }, modifier = Modifier.fillMaxSize()) { mv ->
    mv.getMapAsync { map ->
        map.moveCamera(CameraUpdateFactory.newLatLngZoom(BRATISLAVA, 13f))
    }
}
\end{lstlisting}

\texttt{MapView} disponuje vlastným životným cyklom, ktorý je potrebné manuálne prepojiť s životným cyklom obrazovky pomocou \texttt{DisposableEffect}. Tým sa zabezpečí, že mapa správne reaguje na udalosti ako \textit{onPause}, \textit{onResume} a \textit{onDestroy}.

\subsubsection{GPS poloha a správa oprávnení}
Aplikácia vyžaduje oprávnenie \texttt{ACCESS\_FINE\_LOCATION}, po ktorého udelení získa poslednú známu polohu zariadenia pomocou \texttt{FusedLocationProviderClient}. Kamera mapy sa presunie na polohu používateľa iba raz pri inicializácii. Tento stav je riadený príznakom \texttt{mapMovedToUser} v \texttt{AppViewModel}, čím sa zabraňuje nežiaducemu a opakovanému presunu kamery pri každej rekomponzícii rozhrania.

\begin{lstlisting}[language=Kotlin, caption=Získanie aktuálnej polohy používateľa]
LocationServices.getFusedLocationProviderClient(context)
    .lastLocation
    .addOnSuccessListener { loc ->
        userLocation = LatLng(loc.latitude, loc.longitude)
        appViewModel.updateUserLocation(loc.latitude, loc.longitude)
    }
\end{lstlisting}

\subsubsection{Práca s markermi a interakcia}
Všetky udalosti sú konvertované na objekty \texttt{MarkerOptions} pomocou metódy \texttt{toMarkerOptions()}, ktorá je definovaná priamo v dátovom modeli \texttt{Event}.

\begin{lstlisting}[language=Kotlin, caption=Konverzia modelu Event na MarkerOptions]
fun toMarkerOptions(icon: BitmapDescriptor? = null): MarkerOptions =
    MarkerOptions()
        .position(LatLng(latitude, longitude))
        .title(title)
        .apply { icon?.let { icon(it) } }
\end{lstlisting}

Zobrazované markery sú synchronizované so stavom \texttt{appViewModel.activeEvents} prostredníctvom \texttt{snapshotFlow}. Pri každej zmene zoznamu udalostí sa mapa reaktívne aktualizuje a pridá nové body s vlastným grafickým vizuálom (čierny pin). Do vlastnosti \texttt{tag} každého markera sa ukladá unikátne \texttt{event.id}, čo umožňuje presnú identifikáciu udalosti pri interakcii. Po kliknutí na marker aplikácia identifikuje udalosť a naviguje používateľa na obrazovku s detailom.

\paragraph{Automatické skrývanie skončených udalostí:}
\texttt{AppViewModel} vystavuje computed property \texttt{activeEvents}, ktorá
filtruje interné pole \texttt{\_events} a vracia len udalosti, ktorých koniec
je v budúcnosti:

\begin{lstlisting}[language=Kotlin, caption=Computed property activeEvents v AppViewModel]
val activeEvents: List<Event>
    get() = _events.values.filter { event ->
        val end = event.dateTo?.time ?: event.dateFrom?.time
        end == null || end > System.currentTimeMillis()
    }
\end{lstlisting}

Pravidlá filtrovania sú:
\begin{itemize}
    \item Ak udalosť má \texttt{dateTo}, rozhoduje tento čas.
    \item Ak \texttt{dateTo} neexistuje, použije sa \texttt{dateFrom}.
    \item Udalosti bez dátumu sú vždy zobrazené.
\end{itemize}

Markery na mape aj filter sheet používajú výhradne \texttt{activeEvents} —
skončené udalosti zmiznú z UI automaticky bez potreby manuálneho mazania
z databázy. Interné \texttt{\_events} zostáva kompletné, aby \texttt{checkNoShows()}
stále videl aj minulé eventy pri kontrole neúčastí.

\textit{Dizajnové rozhodnutie:} Udalosti sa fyzicky nemažú z Firestore — zostávajú
pre históriu účasti, no-show kontrolu a prípadné štatistiky. Skrývanie prebieha
čisto na úrovni UI.

\subsubsection{Okruh pôsobnosti a filtrovanie}
Pre lepšiu prehľadnosť aplikácia filtruje udalosti na základe ich vzdialenosti od používateľa. Markery nachádzajúce sa mimo nastaveného okruhu sú automaticky skryté (\texttt{marker.isVisible = false}). Tento okruh je na mape vizualizovaný ako žltý polopriesvitný kruh (\textit{Circle}), ktorého polomer je prepojený s nastaveniami v \texttt{AppViewModel}.

\subsubsection{Personalizácia zobrazenia}
Používateľ má možnosť prepínať medzi piatimi štýlmi mapového podkladu: \textit{NORMAL, SATELLITE, HYBRID, TERRAIN} a špeciálnym \textit{DARK} režimom, ktorý využíva vlastnú definíciu štýlu vo formáte JSON.